\section{SPEAR: Parameter-Aware Self-Improvement}
\label{sec:method}

\begin{figure*}[t]
\centering
\begin{tikzpicture}[>=Latex, every node/.style={font=\small},
  stage/.style={draw=blue!55!black,rounded corners=3pt,thick,
  fill=blue!3,align=center,minimum height=0.95cm,text width=3.1cm},
  minor/.style={draw=gray!65,rounded corners=3pt,
  fill=gray!5,align=center,minimum height=0.85cm,text width=3.1cm}]
\node[stage] (failure) at (1.7,0) {\textbf{Agent failure state}\\trace, checkpoint, budget};
\node[stage] (query) at (5.7,0) {\textbf{Parameter-aware query}\\estimate PI-VOI; choose source};
\node[stage] (evidence) at (9.7,0) {\textbf{Evidence acquisition}\\teacher / retrieval / probe};
\node[stage] (propose) at (13.7,0) {\textbf{Functional intervention}\\basis coefficients + residual};
\node[minor] (verify) at (13.7,-1.85) {\textbf{Paired validation}\\retention, cost, rollback};
\node[minor] (replay) at (5.7,-1.85) {\textbf{Outcome replay}\\train query and update policies};
\draw[->,thick] (failure.east)--(query.west);
\draw[->,thick] (query.east)--(evidence.west);
\draw[->,thick] (evidence.east)--(propose.west);
\draw[->,thick] (propose.south)--(verify.north);
\draw[->,thick] (verify.west)--node[above,font=\scriptsize]{measured intervention outcome}(replay.east);
\draw[->,thick] (replay.north)--(query.south);
\draw[->,dashed,gray] (verify.south) -- ++(0,-0.37) -| (failure.south);
\end{tikzpicture}
\caption{Proposed \spear{} loop. The diagnostic policy is trained on
\emph{verified post-update outcomes}, not the superficial correctness of
an explanation. A candidate can be rejected or replaced by a no-update
decision; only an accepted update persists to future episodes.}
\label{fig:architecture}
\end{figure*}

\subsection{Decision State and the Update Objective}
Let $\theta_t\in\R^d$ denote the \emph{persistent} parameters of a
student agent at improvement round $t$. A long-horizon interaction
produces a partial trajectory $h_t=(o_0,a_0,\ldots,o_k)$ that contains
an observed failure or capability gap. The decision state is
\begin{equation}
s_t=(h_t,\theta_t,D_t,b_t,\M_t),
\end{equation}
where $D_t$ is a functional update basis, $b_t$ is the remaining
resource budget, and $\M_t$ is a log of prior evidence decisions,
parameter proposals, and evaluated outcomes. A source action
$a\in\A_{\mathrm{src}}$ may query a teacher, retrieve external
documentation, run a controlled environment probe, or stop.
An outcome $e\sim P_t(e\mid s_t,a)$ includes the acquired content,
source, cost, timestamp, and provenance. The distribution $P_t$ is
unknown and may vary with model checkpoint and environment version.

The available parameter interventions constitute a bounded set
$\Delta_t$, augmented with $\delta=0$; the latter is important because
the correct decision can be to retain the current checkpoint.
A \emph{memory-only} remedy is a separate non-parametric comparator:
it may improve current-task utility but does not count as durable
parameter improvement.

Let $\mathcal{T}^{\rm new}_t$ be independent future task instances from
the target skill distribution and $\mathcal{T}^{\rm old}_t$ a retention
suite of previously learned tasks. Denote strict episode success by
$S(\theta;z)\in[0,1]$, evaluated by the environment's state checker.
Define the net \emph{update value}
\begin{align}
J_t(\delta) ={}&
  \E_{z\sim\mathcal{T}^{\rm new}_t}[S(\theta_t+\delta;z)]
     \nonumber\\
 &-\lambda_F\Big[
     \E_{z\sim\mathcal{T}^{\rm old}_t}
     \big(S(\theta_t;z)-S(\theta_t+\delta;z)\big)
   \Big]_+ \nonumber\\
 &-\lambda_U C_{\rm update}(\delta).
\label{eq:updatevalue}
\end{align}
Here $[x]_+=\max(x,0)$ and $C_{\rm update}$ accounts for actual
training, evaluation, and deployment cost; the separate source cost is
denoted $c_{\rm src}(a)$. The weights and normalization must be fixed
\emph{before} testing. Performance on a final held-out test set must
not enter training, policy selection, or the acceptance gate.

\subsection{Parameter-Intervention Value of Information}
\label{sec:pivoi}
Given history $\M_t$, define the conditional expected utility of
an intervention after observing possible evidence $e$:
\begin{equation}
\mu_t(\delta\mid s_t,e,\M_t)
=\E[J_t(\delta)\mid s_t,e,\M_t].
\label{eq:mu}
\end{equation}
The proposed one-step \emph{parameter-intervention value of
information} for a source action is
\begin{align}
V^{\rm PI}_t(a\mid s_t,\M_t)
={}&\E_{e\sim P_t(\cdot\mid s_t,a)}
  \left[
  \max_{\delta\in\Delta_t\cup\{0\}}
  \mu_t(\delta\mid s_t,e,\M_t)
  \right]\nonumber\\
 &-\max_{\delta\in\Delta_t\cup\{0\}}
  \mu_t(\delta\mid s_t,\M_t)
  -c_{\rm src}(a).
\label{eq:pivoi}
\end{align}
The second term is the best available parameter decision without
paying for further evidence. Crucially, a source is not rewarded
for reducing uncertainty unless the reduction improves the
\emph{downstream intervention choice}. This is a decision-dependent
instance of classical value of information, \emph{not} a new
information-theoretic identity.

\begin{proposition}[Nonnegative gross value under ideal decisions]
\label{prop:voi}
Suppose the source does not itself alter $\theta_t$, all candidate
interventions remain available, the utility model is Bayes-correct,
and source cost is zero. Then the first two terms of
Eq.~\eqref{eq:pivoi} have nonnegative difference.
\end{proposition}
\noindent\textit{Proof.}
Let $\delta^\star$ maximize expected utility before acquisition.
After seeing $e$, the decision-maker can always select
$\delta^\star$. Taking expectations over $e$ and applying the law
of iterated expectations yields an expected post-query optimum
no smaller than the pre-query optimum. \hfill$\square$

This result is only an ideal benchmark: when queries are costly,
evidence is distribution-shifting, or $\mu_t$ is miscalibrated,
the estimated net value and the \emph{realized} gain can be negative.
A myopic policy chooses
$a_t=\arg\max_a\widehat V^{\rm PI}_t(a\mid s_t,\M_t)$
when that maximum is positive, otherwise stopping.
Sequential acquisition re-evaluates the same criterion after each
observation within budget. A one-step approximation does not
generally solve the optimal multi-step experimental-design problem.

\subsection{Estimating Parameter-Aware Evidence Value}
\label{sec:valuelearning}
We propose a learned intervention-value model
$Q_\omega(s,e,\delta)$, trained to predict the measured change
$J_t(\delta)-J_t(0)$, together with a calibrated uncertainty estimate
$\sigma_\omega$. For each randomized evidence decision and feasible
candidate $\delta$, the training log stores the \emph{actual} paired
validation outcome, cost, environment seed, and versioned checkpoint.
The conditional evidence model
$\widehat P_\xi(e\mid s,a)$ is optional; a direct
source-value regressor is a lower-complexity alternative when sampling
full textual evidence is impractical.

A practical source-value estimator simulates a limited set of
candidate evidence outcomes and updates:
\begin{equation}
\widehat V^{\rm PI}(a)
= \frac{1}{K}\sum_{i=1}^K
 \max_{\delta\in\widehat\Delta(s,e_i)}
 Q_\omega(s,e_i,\delta)
 -\max_{\delta\in\widehat\Delta(s)}
 Q_\omega(s,\varnothing,\delta)-c_{\rm src}(a).
\label{eq:est}
\end{equation}
Here $e_i\sim\widehat P_\xi(\cdot\mid s,a)$ and every candidate set
includes zero. This is an \emph{approximation}, not an unbiased
counterfactual estimator. Limited candidate search and maximization
of noisy scores can be optimistic; calibration on a separate
validation split, uncertainty penalties, and occasional random
source exploration are necessary. Logged source propensities must
be retained for subsequent off-policy evaluation.

\subsection{Evidence-to-Functional Intervention}
\label{sec:generator}
Let $D_t\in\R^{d\times k}$ comprise $k\ll d$ reusable parameter
directions from previously \emph{validated} updates. The proposed
generator predicts both basis coefficients and a bounded residual:
\begin{equation}
\delta_\phi(s_t,e)=
    D_t g_\phi(s_t,e)+r_\phi(s_t,e),\qquad
    \|r_\phi\|_2\le \kappa,\ \|\delta_\phi\|_2\le R.
\label{eq:functional}
\end{equation}
For large weight matrices, $r_\phi$ can be represented by small
low-rank factors or omitted in the minimal implementation.
The generator receives structured history and evidence provenance,
not just the consultant's natural-language answer. A finite library
of pre-trained directions with learned mixture coefficients is a
crucial low-cost baseline for the full generator.

A value-grounded training objective is
\begin{align}
\mathcal{L}_{\rm int}(\phi)
={}&-\E_{(s,e)\sim\M}
 \big[Q_\omega(s,e,\delta_\phi(s,e))\big]\nonumber\\
 &+\lambda_{\rm res}\|r_\phi(s,e)\|_2^2
 +\lambda_{\rm trust}\|\delta_\phi(s,e)\|_2^2.
\label{eq:generatorloss}
\end{align}
The model must not trust this surrogate alone; it proposes candidates
for \emph{real execution}, and a high predicted value is not a
validation pass. Direction reuse can be expanded only after
independently verified improvements. Directions must be attached
to a compatible checkpoint, layer layout, and parameter
normalization; cross-backbone transfer is not assumed.

\subsection{Counterfactual Acceptance and Persistence}
\label{sec:verification}
Each candidate update is evaluated against the \emph{same}
pre-update checkpoint on paired, independently initialized task
instances. For every evaluation seed, run both $\theta_t$ and
$\theta_t+\delta$ from equivalent environment state, recording
strict success and side effects. Keep data used to create the
candidate disjoint from its validation data, and never select on
the final test split. An acceptance rule can require a positive
pre-specified lower confidence bound on the new-skill gain and an
old-skill regression below tolerance $\epsilon_F$:
\begin{equation}
\mathrm{LCB}_{1-\alpha}(\Delta U_{\rm new})>0,\quad
\mathrm{UCB}_{1-\alpha}(F_{\rm old})\le\epsilon_F.
\label{eq:accept}
\end{equation}
The confidence construction, multiplicity correction, and number
of repeats must be pre-registered. When either check fails,
rollback to $\theta_t$; accept $\delta=0$ if no proposal passes.

When accepted, \emph{persist} the update as a versioned adapter or
merge it into a new checkpoint, yielding
$\theta_{t+1}=\theta_t+\delta_t^\star$.
At the next evaluation episode, disable the original teacher
message and retrieved document. If improvements require reloading
temporary query-specific adapters or re-injecting the diagnosis,
classify the effect as transient adaptation rather than persistent
parameter learning.

\subsection{Recursive Learning of the Update Mechanism}
\label{sec:recursive}
Define the learned self-improvement machinery as
$\psi_t=(\eta_t,\phi_t,\omega_t,\xi_t)$, where $\eta_t$
is the diagnostic source policy, $\phi_t$ the generator, and the
remaining components are value/evidence predictors.
After each verified trial, update
\begin{equation}
(\theta_t,\psi_t,\M_t)
\longrightarrow
(\theta_{t+1},\psi_{t+1},\M_{t+1}).
\label{eq:recursion}
\end{equation}
A constrained contextual-bandit or off-policy policy-gradient
loss for $\eta$ can use \emph{measured intervention gain minus
acquisition cost}, with recorded randomized action propensities.
The generator is trained with Eq.~\eqref{eq:generatorloss} and
ground-truth acceptance/rejection feedback. All outcomes are
versioned because $Q_\omega$ at a prior $\theta_t$ need not
generalize to the next checkpoint.

\begin{algorithm}[t]
\caption{Proposed SPEAR improvement round (research protocol)}
\label{alg:spear}
\DontPrintSemicolon
\KwIn{checkpoint $\theta_t$, failure history $h_t$, budget $b_t$,
      replay $\M_t$, candidate set $\Delta_t\cup\{0\}$}
initialize $e\leftarrow\varnothing$; source log $L\leftarrow\varnothing$\;
\While{budget remains}{
  estimate source values $\widehat V^{\rm PI}(a\mid s_t,e,\M_t)$\;
  \If{$\max_a\widehat V^{\rm PI}(a)\le 0$}{break}
  choose source $a$ (with logged exploration propensity)\;
  acquire evidence $e_a$; append $(a,e_a,c_a)$ to $L$\;
  update evidence context $e\leftarrow e\cup\{e_a\}$\;
}
generate bounded candidate updates from $(s_t,e)$, including $0$\;
evaluate candidates and unchanged checkpoint on paired validation\;
accept guarded best candidate or rollback to $\theta_t$\;
log source actions, all candidate outcomes, and retention checks\;
update $\eta,\phi,\omega$ from intervention-outcome replay\;
\KwOut{persistent $\theta_{t+1}$, updated mechanism, versioned $\M_{t+1}$}
\end{algorithm}

A claim of \emph{recursive self-improvement} must be based on an
\emph{incremental effect of learning $\psi$}. Growth of $\M_t$ alone
is not sufficient: a frozen source policy and frozen generator
must be compared with the adaptive versions under \emph{the same}
memory content, number of candidate evaluations, and task stream.
Only a reliable positive meta-learning effect supports the
recursive part of the claim.

\subsection{Distinguishing Information Gain From Rephrasing}
If a teacher diagnosis is a deterministic function of the entire
observed trace, it adds no Shannon information about a latent
environment state conditional on that trace. It can nonetheless
provide a computational advantage to a bounded-capacity generator
by making the relevant pattern easier to use. The evaluation
therefore separates (i) \emph{new external observations} from
(ii) \emph{trace-only linguistic transformations}, and does not
conflate either with successful weight-space learning.
